% preamble.tex -- mise en forme du rapport de stage ENSIASD
% Times New Roman, 12 pt, interligne 1,5, marges 2,5 cm,
% titres de chapitre 16 pt MAJUSCULES gras alignes a droite.
% Palette : bleu clair (accent principal) + rouge sang fonce (accent
% secondaire, discret) sur fond blanc -- sobre, dans l'esprit des
% rapports techniques Apple / OpenAI.

\usepackage{iftex}
\ifPDFTeX
  \usepackage[T1]{fontenc}
  \usepackage[utf8]{inputenc}
  \usepackage{mathptmx}
  \usepackage{courier}
\else
  \usepackage{fontspec}
  \defaultfontfeatures{Ligatures=TeX}
  \setmainfont{Times New Roman}
  \setmonofont{Courier New}[Scale=0.86]
\fi

\usepackage[french]{babel}
\usepackage{csquotes}
\usepackage[
  a4paper,
  top=2.5cm,
  bottom=2.5cm,
  left=2.5cm,
  right=2.5cm
]{geometry}
\usepackage{setspace}
\onehalfspacing
\usepackage{microtype}
\usepackage{xcolor}
\usepackage{graphicx}
\usepackage{tikz}
\usetikzlibrary{arrows.meta,positioning,fit,calc,shapes.geometric,shadows.blur,backgrounds}
\usepackage{booktabs}
\usepackage{array}
\usepackage{tabularx}
\usepackage{colortbl}
\usepackage{multirow}
\usepackage{makecell}
\usepackage{caption}
\usepackage{subcaption}
\usepackage{titlesec}
\usepackage{tocloft}
\usepackage{fancyhdr}
\usepackage{enumitem}
\usepackage{listings}
\usepackage{float}
\usepackage{needspace}
\usepackage{etoolbox}
\usepackage{ragged2e}
\usepackage{xurl}
\usepackage{pgfgantt}
\usepackage{hyperref}

% ---- palette : bleu clair + rouge sang (discret) sur blanc -----------------
\definecolor{accentblue}{HTML}{2E6DA4}      % accent principal -- structure, liens
\definecolor{bluedeep}{HTML}{1B4A73}        % titres de chapitre, texte fort en bleu
\definecolor{bluesoft}{HTML}{AFC9E0}        % filets fins, fonds tres legers
\definecolor{bluewash}{HTML}{EEF3F8}        % fond d'encadre "information"
\definecolor{bloodred}{HTML}{7A1F2B}        % accent secondaire -- alerte, limite
\definecolor{bloodredsoft}{HTML}{E7CACD}    % filet / fond tres leger, alerte
\definecolor{bloodredwash}{HTML}{FBF0F1}    % fond d'encadre "limite / attention"
\definecolor{ink}{HTML}{1B1B1D}             % texte courant, quasi noir
\definecolor{muted}{HTML}{68686C}           % texte secondaire (pied de page)
\definecolor{rulegray}{HTML}{DADADD}        % filets de tableaux
\definecolor{codebg}{HTML}{F5F5F7}          % fond des extraits de code

\color{ink}

% ---- liens -------------------------------------------------------------------
\hypersetup{
  colorlinks=true,
  linkcolor=bluedeep,
  citecolor=bluedeep,
  urlcolor=bluedeep,
  bookmarksnumbered=true,
  pdftitle={Rapport de stage -- Orchestration d'outils d'analyse de securite statique},
  pdfauthor={\AuthorName}
}

% ---- legendes : figures en dessous, tableaux au-dessus, 11 pt --------------
\captionsetup{
  font={small},
  labelfont={it},
  textfont={it},
  justification=centering,
  skip=6pt
}
\captionsetup[figure]{position=bottom, font={small,it}, labelfont={it}}
\captionsetup[table]{position=top, font={small}, labelfont={bf}, textfont={}}
\renewcommand{\thefigure}{\arabic{chapter}.\arabic{figure}}
\renewcommand{\thetable}{\arabic{chapter}.\arabic{table}}

% ---- titres de chapitre (ENSIASD) ------------------------------------------
% Nouvelle page, 16 pt, MAJUSCULES, gras, aligne a droite,
% espacement avant 0 / apres 54 pt. Un filet bleu sobre en dessous.
\titleformat{\chapter}[display]
  {\normalfont\bfseries\color{bluedeep}\fontsize{16pt}{20pt}\selectfont}
  {\filleft\color{accentblue}\fontsize{11pt}{14pt}\selectfont CHAPITRE \thechapter}
  {8pt}
  {\filleft\MakeUppercase}
  [\vspace{8pt}{\color{accentblue}\titlerule[2.2pt]}\vspace{3pt}{\color{bloodredsoft}\titlerule[0.5pt]}]
\titlespacing*{\chapter}{0pt}{0pt}{54pt}

\titleformat{name=\chapter,numberless}[display]
  {\normalfont\bfseries\color{bluedeep}\fontsize{16pt}{20pt}\selectfont}
  {}
  {0pt}
  {\filleft\MakeUppercase}
  [\vspace{8pt}{\color{accentblue}\titlerule[2.2pt]}\vspace{3pt}{\color{bloodredsoft}\titlerule[0.5pt]}]
\titlespacing*{name=\chapter,numberless}{0pt}{0pt}{54pt}

% Titres de paragraphe : I, II, III -- 14 pt gras, avant/apres 12 pt
\renewcommand{\thesection}{\Roman{section}}
\renewcommand{\thesubsection}{\thesection.\arabic{subsection}}
\renewcommand{\thesubsubsection}{\thesubsection.\arabic{subsubsection}}

\titleformat{\section}
  {\normalfont\bfseries\fontsize{14pt}{18pt}\selectfont\color{ink}}
  {\thesection.}{0.75em}{}
\titlespacing*{\section}{0pt}{12pt}{12pt}

\titleformat{\subsection}
  {\normalfont\bfseries\fontsize{12pt}{16pt}\selectfont\color{ink}}
  {\thesubsection.}{0.6em}{}
\titlespacing*{\subsection}{0pt}{6pt}{6pt}

\titleformat{\subsubsection}
  {\normalfont\bfseries\fontsize{12pt}{16pt}\selectfont\color{ink}}
  {\thesubsubsection.}{0.6em}{}
\titlespacing*{\subsubsection}{0pt}{6pt}{6pt}

% ---- table des matieres / listes -------------------------------------------
\renewcommand{\cftchapfont}{\bfseries\color{bluedeep}}
\renewcommand{\cftchappagefont}{\bfseries\color{bluedeep}}
\renewcommand{\cftsecfont}{\color{ink}}
\renewcommand{\cftsecpagefont}{\color{ink}}
\renewcommand{\cftdot}{.}
\renewcommand{\cftdotsep}{4.5}
\setlength{\cftbeforechapskip}{8pt}
\renewcommand{\cftchapleader}{\cftdotfill{\cftdotsep}}

% ---- en-tetes / pieds de page : numero en bas a droite ---------------------
\pagestyle{fancy}
\fancyhf{}
\renewcommand{\headrulewidth}{0pt}
\renewcommand{\footrulewidth}{0.6pt}
\renewcommand{\footrule}{{\color{bluesoft}\hrule width\headwidth height 0.6pt}}
\fancyfoot[L]{\footnotesize\color{muted}\AuthorName}
\fancyfoot[R]{\bfseries\color{bluedeep}\thepage}
\fancypagestyle{plain}{
  \fancyhf{}
  \renewcommand{\headrulewidth}{0pt}
  \renewcommand{\footrulewidth}{0.6pt}
  \renewcommand{\footrule}{{\color{bluesoft}\hrule width\headwidth height 0.6pt}}
  \fancyfoot[L]{\footnotesize\color{muted}\AuthorName}
  \fancyfoot[R]{\bfseries\color{bluedeep}\thepage}
}

% Pages liminaires : pas de numerotation arabe visible
\fancypagestyle{front}{
  \fancyhf{}
  \renewcommand{\headrulewidth}{0pt}
  \renewcommand{\footrulewidth}{0pt}
  \fancyfoot[R]{\color{bluesoft}\thepage}
}

% ---- listes -------------------------------------------------------------------
\setlist[itemize]{leftmargin=1.4em, itemsep=0.25em, topsep=0.4em}
\setlist[enumerate]{leftmargin=1.6em, itemsep=0.25em, topsep=0.4em}

% ---- extraits de code ----------------------------------------------------------
\lstdefinestyle{code}{
  basicstyle=\ttfamily\fontsize{10pt}{13pt}\selectfont\color{ink},
  backgroundcolor=\color{codebg},
  frame=tb,
  rulecolor=\color{accentblue},
  framesep=8pt,
  xleftmargin=8pt,
  xrightmargin=8pt,
  breaklines=true,
  breakatwhitespace=true,
  columns=fullflexible,
  keepspaces=true,
  showstringspaces=false,
  commentstyle=\color{muted}\itshape,
  keywordstyle=\bfseries\color{bluedeep},
  stringstyle=\color{bloodred},
  morekeywords={class,def,return,from,import,if,else,for,in,None,True,False,self},
}
\lstset{style=code, gobble=0, captionpos=b}

% The listings package scans its own captions/options character by
% character under pdfLaTeX, and can choke on raw UTF-8 accented bytes
% coming from inputenc (error: "Invalid UTF-8 byte sequence (\lst@EC)").
% This table tells listings how to redraw every accented character our
% French captions use, so it never has to guess.
\lstset{
  extendedchars=true,
  literate=
    {á}{{\'a}}1 {é}{{\'e}}1 {í}{{\'i}}1 {ó}{{\'o}}1 {ú}{{\'u}}1
    {à}{{\`a}}1 {è}{{\`e}}1 {ì}{{\`i}}1 {ò}{{\`o}}1 {ù}{{\`u}}1
    {â}{{\^a}}1 {ê}{{\^e}}1 {î}{{\^i}}1 {ô}{{\^o}}1 {û}{{\^u}}1
    {ä}{{\"a}}1 {ë}{{\"e}}1 {ï}{{\"i}}1 {ö}{{\"o}}1 {ü}{{\"u}}1
    {ç}{{\c c}}1 {œ}{{\oe}}1
    {É}{{\'E}}1 {È}{{\`E}}1 {Ê}{{\^E}}1 {À}{{\`A}}1 {Ç}{{\c C}}1
}

% ---- tableaux : en-tete bleu ---------------------------------------------------
\newcommand{\tablehead}{\rowcolor{bluedeep}\color{white}}
\newcolumntype{L}{>{\raggedright\arraybackslash}X}
\newcolumntype{C}{>{\centering\arraybackslash}X}
\arrayrulecolor{rulegray}

% ---- encadres (information en bleu, limite/alerte en rouge sang) -----------
\newenvironment{infobox}[1]{%
  \par\needspace{4\baselineskip}
  \begin{center}
  \begin{tabularx}{\linewidth}{!{\color{accentblue}\vrule width 2.4pt}p{0.4cm}X}
    & & \\[-0.5em]
    & \multicolumn{1}{l}{\bfseries\color{bluedeep} #1} \\
    & &
  \end{tabularx}
  \end{center}
  \vspace{-1.6em}
  \begin{tabularx}{\linewidth}{!{\color{accentblue}\vrule width 2.4pt}p{0.4cm}X}
    & &
}{%
  \end{tabularx}
  \vspace{0.4em}
}

\newenvironment{warnbox}[1]{%
  \par\needspace{4\baselineskip}
  \begin{center}
  \begin{tabularx}{\linewidth}{!{\color{bloodred}\vrule width 2.4pt}p{0.4cm}X}
    & & \\[-0.5em]
    & \multicolumn{1}{l}{\bfseries\color{bloodred} #1} \\
    & &
  \end{tabularx}
  \end{center}
  \vspace{-1.6em}
  \begin{tabularx}{\linewidth}{!{\color{bloodred}\vrule width 2.4pt}p{0.4cm}X}
    & &
}{%
  \end{tabularx}
  \vspace{0.4em}
}

% ---- utilitaires ----------------------------------------------------------------
\newcommand{\tool}[1]{\texttt{#1}}
\newcommand{\modfile}[1]{\texttt{#1}}

\newcommand{\frontchapter}[1]{%
  \clearpage
  \thispagestyle{front}%
  \chapter*{#1}%
  \addcontentsline{toc}{chapter}{#1}%
}

% Corps du texte : justifie, Times 12 pt, interligne 1,5 (ENSIASD).
\justifying
\setlength{\parindent}{1.25em}
\setlength{\parskip}{0.35em}

\renewcommand{\contentsname}{TABLE DES MATI\`ERES}
\renewcommand{\listfigurename}{LISTE DES FIGURES}
\renewcommand{\listtablename}{LISTE DES TABLEAUX}
\renewcommand{\bibname}{R\'EF\'ERENCES BIBLIOGRAPHIQUES}
